% =====================================================================
% Kapitel 9: Tests
% =====================================================================

\chapter{Tests}

\label{ch:tests}

Tests sind das Sicherheitsnetz des Projekts. Sie prüfen, ob der
Code das tut, was er tun soll -- und verhindern, dass spätere
Änderungen bestehende Funktionalität brechen.

\section{Test-Philosophie}

\begin{merke}
Das Projekt folgt der \emph{Pyramide} der Tests:

\begin{itemize}
  \item \textbf{Viele} schnelle \emph{Unit-Tests} (einzelne Funktionen).
  \item \textbf{Wenige} \emph{Integration-Tests} (mehrere
        Komponenten zusammen).
  \item \textbf{Manuelle} \emph{Checks} für Dinge, die schwer zu
        automatisieren sind (visuelle TUI-Korrektheit).
\end{itemize}
\end{merke}

\section{Python-Tests mit pytest}

\subsection{Überblick}

\begin{lstlisting}[language=bash, caption={Python-Tests laufen}]
.venv/bin/pytest
# 55 passed
\end{lstlisting}

55~Tests sind in der Suite (Stand: Handbuch-Erstellung). Sie sind
in folgende Module aufgeteilt:

\begin{itemize}
  \item \code{tests/test\_config.py} -- Konfigurationsladen.
  \item \code{tests/test\_models.py} -- SQLAlchemy-ORM-Registrierung.
  \item \code{tests/test\_fred.py} -- FRED-Vendor (mit gemocktem
        HTTP).
  \item \code{tests/test\_ecb.py} -- ECB-Vendor (mit gemocktem HTTP).
  \item \code{tests/test\_yahoo.py} -- Yahoo-Vendor (mit
        gemocktem yfinance).
  \item \code{tests/analytics/test\_metrics.py} -- Returns + Metriken.
\end{itemize}

\subsection{Test-Fixtures}

\begin{definition}
\textbf{pytest Fixture}
Eine \emph{Fixture} ist eine wiederverwendbare Test-Setup-Funktion.
Sie wird mit \code{@pytest.fixture} dekoriert und kann in Tests
\emph{injiziert} werden, indem sie als Argument mit dem gleichen
Namen wie die Fixture-Funktion deklariert wird.
\end{definition}

\begin{lstlisting}[language=python, caption={pytest-Fixture für Settings-Cache-Reset},label=lst:fixture]
@pytest.fixture(autouse=True)
def _reset_settings_cache():
    """Reset the pydantic-settings lru_cache around each test."""
    get_settings.cache_clear()
    yield
    get_settings.cache_clear()
\end{lstlisting}

Diese \code{autouse=True}-Fixture wird \emph{vor jedem Test}
automatisch ausgeführt und stellt sicher, dass der
\code{lru\_cache} der \code{get\_settings()}-Funktion geleert
ist -- sonst würden Test-Pollution und unstabile Ergebnisse
entstehen.

\subsection{Mocking}

\begin{definition}
\textbf{Mock}
Ein \emph{Mock} ist ein gefälschtes Objekt, das das Verhalten eines
realen Objekts nachahmt. In Tests verwenden wir Mocks, um externe
Abhängigkeiten (HTTP-Server, Datenbanken) zu simulieren, ohne
\emph{tatsächlich} mit ihnen zu sprechen -- das macht Tests
\emph{ schnell}, \emph{ deterministisch } und \emph{ offline
  lauffähig}.
\end{definition}

Listing~\ref{lst:mock-fred} zeigt einen typischen Mock-Test für
den FRED-Vendor.

\begin{lstlisting}[language=python, caption={FRED-Vendor mit gemocktem HTTP},label=lst:mock-fred]
def test_fred_fetch_parses_response(fake_fred_observations, monkeypatch):
    monkeypatch.setenv("FRED_API_KEY", "test-key")
    v = FredVendor()
    v._get = MagicMock(return_value=fake_fred_observations)
    df = v.fetch("DGS3MO", date(2024, 6, 1), date(2024, 6, 30))
    assert df.height == 3
    assert df["value"][0] == pytest.approx(5.02)
    v._get.assert_called_once()
\end{lstlisting}

Die Strategie: \code{v.\_get} wird durch einen \code{MagicMock}
ersetzt, der \emph{jeden} Aufruf mit dem vorbereiteten Fake-Response
beantwortet. So können wir die \emph{Response-Parsing}-Logik
isoliert testen, ohne einen echten HTTP-Server zu kontaktieren.

\section{TypeScript-Tests mit bun test}

\subsection{Überblick}

\begin{lstlisting}[language=bash, caption={TUI-Tests laufen}]
cd tui && bun test
# 27 pass, 0 fail
\end{lstlisting}

\subsection{Mocking in bun test}

Bun's Test-Runner hat eine Jest-kompatible API:
\code{describe}, \code{test}, \code{expect}, \code{spyOn}, \code{mock}.

\begin{lstlisting}[language=typescript, caption={Mocking der pg-Layer im TUI},label=lst:bun-mock]
import { describe, it, expect, spyOn, type Mock } from "bun:test";
import * as dbModule from "./db";
import { listSeries } from "./queries";

let spies: Mock<typeof dbModule.query>[] = [];

function mockQuery<T>(rows: T[]): Mock<typeof dbModule.query> {
  const spy = spyOn(dbModule, "query") as Mock<typeof dbModule.query>;
  spy.mockReset();
  spy.mockResolvedValue(rows as never);
  spies.push(spy);
  return spy;
}

afterEach(() => {
  for (const s of spies) s.mockRestore();
  spies = [];
});

describe("listSeries", () => {
  it("calls query() and returns the rows", async () => {
    const rows: SeriesRow[] = [{...}];
    const spy = mockQuery(rows);
    const out = await listSeries();
    expect(out).toEqual(rows);
    expect(spy).toHaveBeenCalledTimes(1);
    const sql = spy.mock.calls[0]?.[0] as string;
    expect(sql).toContain("UNION ALL");
  });
});
\end{lstlisting}

Wichtige Pattern:

\begin{itemize}
  \item \code{spyOn} ersetzt eine Methode durch einen Spy, der
        Aufrufe aufzeichnet.
  \item \code{mockResolvedValue} setzt den Rückgabewert.
  \item \code{afterEach} restauriert die Spies, um Test-Pollution zu
        vermeiden.
\end{itemize}

\subsection{Integration-Tests}

\begin{lstlisting}[language=typescript, caption={Integration-Test mit echtem Postgres}]
beforeAll(async () => {
  try {
    await query("SELECT 1::int AS x");
    dbReachable = true;
  } catch (e) {
    console.warn("Postgres unreachable, tests will no-op.");
  }
});

function guarded(name, fn) {
  return async () => {
    if (!dbReachable) {
      console.warn(`skip ${name}: DB unreachable`);
      return;
    }
    await fn();
  };
}

it("equity rows carry the YAHOO ticker (AAPL/MSFT/^GSPC)",
   guarded("equity-tickers", async () => {
     const rows = await listSeries();
     const tickers = rows
       .filter(r => r.kind === "equity")
       .map(r => r.id)
       .sort();
     expect(tickers).toEqual(["AAPL", "MSFT", "^GSPC"]);
   }));
\end{lstlisting}

Integration-Tests prüfen das Verhalten gegen die echte
PostgreSQL-DB. Sie sind \emph{no-op} (nicht \emph{fail}), wenn
die DB nicht erreichbar ist -- so bleibt die CI grün, auch wenn
Postgres temporär down ist.

\section{Manuelle Test-Checklist}

Manche Dinge sind schwer zu automatisieren -- visuelle Korrektheit
einer TUI, Tastatur-Feel, Verhalten unter ungewöhnlichen
Bedingungen. Dafür gibt es die manuelle Checkliste in
\code{tui/README.md}.

\begin{beispiel}
\textbf{Auszug aus der manuellen TUI-Checklist}
  \begin{enumerate}
    \item \code{make tui} starten
    \item Header zeigt "12 series" (nicht "loading...")
    \item Erste Zeile (\code{AAPL}) ist selektiert ($\triangleright$)
    \item \code{$\downarrow$} 4x drücken: Selection wandert
    \item \code{r} drücken: Series-Liste wird neu geladen
    \item \code{q} drücken: TUI beendet sauber
  \end{enumerate}
\end{beispiel}

\section{Code-Qualität}

\begin{definition}
\textbf{Linter}
Ein \emph{Linter} ist ein statisches Code-Analyse-Tool, das nach
\emph{Potentiellen Fehlern}, \emph{Stil-Verstößen} und
\emph{Sicherheitsproblemen} sucht, ohne den Code auszuführen.
\end{definition}

Dieses Projekt verwendet:

\begin{itemize}
  \item \textbf{ruff} -- ein extrem schneller Python-Linter
        (in Rust geschrieben). Ersetzt \code{flake8}, \code{isort},
        \code{black} und Teile von \code{pylint}.
  \item \textbf{tsc} -- der TypeScript-Compiler im
        \code{--noEmit}-Modus (nur Typ-Check, keine Ausgabe).
\end{itemize}

\begin{lstlisting}[language=bash, caption={Linter ausführen}]
# Python
.venv/bin/ruff check scripts/ src/ tests/
.venv/bin/ruff format --check scripts/ src/ tests/

# TypeScript
cd tui && bun run typecheck
\end{lstlisting}

\begin{warnung}
Beachte: Das Projekt verwendet \emph{keine} Typen-Annotationen für
\`mypy\` -- die Typen sind in \code{pyproject.toml} unter
\code{[tool.mypy]} deaktiviert. Dies ist eine bewusste
Vereinfachung für ein Retail-Projekt; für ein
\emph{Produktions}-System würde ich strikte Typen empfehlen.
\end{warnung}

\section{Was als Nächstes?}

Mit Tests und Linting als Sicherheitsnetz sind wir bereit für das
\emph{letzte und ambitionierteste} Stück: die
AEGIS-Paper-Replikation (Kapitel~\ref{ch:analytics}). Dort
verbinden wir alle bisherigen Komponenten -- Vendor-Adapter,
ETL-Pipeline, Datenmodell -- zu einer kohärenten quantitativen
Analyse.
